\chapter{保险与风险分担机构}
\label{chap:Insurance_and_Risk_Sharing_Institutions}


\begin{openingcase}{小额保险费与大额承诺}

一个家庭拥有一套价值约200万人民币的房子。虽然在特定年份发生严重火灾的可能性不大，但如果真的发生了火灾，造成的经济损失可能会非常惨重。为了避免承担这种风险，这个家庭购买了保险，并支付每年仅相当于保险公司最终可能需要支付的保险费的一小部分。

这笔交易看起来是不对称的。家庭今天支付的是一笔相对较小且确定的金额，而保险公司则承担潜在的巨大潜在负债。然而，保险公司同时会签订成千上万这样的合同。

\end{openingcase}

\begin{tcolorbox}[
  enhanced,
  colback=FinanceBlue!6,
  colframe=FinanceBlue,
  boxrule=0.6pt,
  sharp corners,
  left=8pt,
  right=8pt,
  top=7pt,
  bottom=7pt,
  before skip=1em,
  after skip=1.6em,
  title=\textbf{核心问题},
  coltitle=white,
  colbacktitle=FinanceBlue,
  fonttitle=\bfseries,
  attach boxed title to top left={xshift=6pt,yshift=-2mm},
  boxed title style={
    boxrule=0pt,
    sharp corners,
    left=5pt,
    right=5pt,
    top=2pt,
    bottom=2pt
  }
]
\begin{enumerate}[
  leftmargin=*,
  label=\textcolor{FinanceBlue}{\textbf{\arabic*.}},
  itemsep=0.55em,
  topsep=0.2em
]


\item 为什么保险能够改善风险分担，以及何时通过集中管理使得整体损失更加可预测？

\item 保险契约如何在解决逆向选择和道德风险的同时，还能保持投保人的积极性？

\item 保险公司如何为保险契约定价、评估和管理其负债，如何管理资产和资本，以及如何通过再保险和资本市场来转移风险？

\end{enumerate}
\end{tcolorbox}

\section{保险、风险汇集与可保性}

保险公司通过合同来转移不确定事件的财务后果，用已知的保险费换取对特定损失的保障。通过结合许多这样的合同，保险公司可以汇集风险，使总损失更加可预测。而这种安排的可行性取决于潜在风险是否可以被保险所涵盖。

\subsection{个体风险与保险需求}

个人可能会面临一种低概率但相对较大的损失，相对于初始财富而言。设$W$表示初始财富，假设随机损失$\tilde{L}$以概率$p$取值$L$，而以概率$1-p$则取值为零。如果没有保险，最终的财富状况……

\begin{equation}
W_T=
\begin{cases}
W-L, & \text{以一定的概率 } p,\\
W, & \text{以一定的概率 } 1-p.
\end{cases}
\label{eq:ch15-uninsured-wealth}
\end{equation}

预期的经济损失为

\begin{equation}
\mathbb{E}[\tilde{L}]=pL.
\label{eq:ch15-expected-loss-simple}
\end{equation}

仅有预期损失，无法概括投保人在不同状态下的风险暴露。保险重新分配各状态下的财富：无论是否出险，投保人均支付保险费；损失发生时则获得赔偿。效用函数对财富呈凹性时，减少状态间财富差异可能提高预期效用。因此，风险厌恶的投保人可能愿意支付超过预期损失的保险费，以换取对重大不利结果的保障。

假设保险费为$P$时，能够获得全额覆盖。忽略排除条款和索赔的障碍后，与这种损失相关的最终财富就成为了

\begin{equation}
W_T^{I}=W-P.
\label{eq:ch15-fully-insured-wealth}
\end{equation}

完全保险以更稳定的财富状况替代状态依存财富。即使保险费超过预期损失，只要减少财富风险的价值超过保费附加，保险仍可能有价值。

保险还可以减少为了应对低频但损失严重的事件而需要持有的预防性流动资产的需求。当在发生损失后借款或清算资产变得成本高昂或极其困难时，这种效应就更加明显了。

\begin{example}{相对于家庭财富的重大损失}{household-insurance-demand}

一个家庭初始财富为100万人民币，全年面临2\%的概率会遭受40万人民币的损失。如果不进行保险处理，那么如果发生损失，最终财富将仅为60万人民币；而如果没有发生损失，最终财富则仍为100万人民币。

预期损失是

$$
0.02\times 400{,}000=8{,}000.
$$

在保险费为$P$的情况下，无论处于哪种状态，最终的财富都约为人民币$1{,}000{,}000-P$。因此，一个厌恶风险的家庭可能会接受高于人民币8,000的保险费，只要减少面临重大不利状况的风险所带来的价值足以覆盖保险费的增加。

\end{example}


\subsection{风险汇集与大数定律}

当保险公司将许多风险汇集在一起时，风险汇集的效果会更好。作为参考，假设有$n$个投保人，他们具有独立的损失指标$X_i$，其中……

$$
X_i=
\begin{cases}
1, & \text{如果投保人 } i \text{ 感受到了损失},\\
0, & \text{否则，请保留LaTeX命令、公式、交叉引用以及数字的内容。},
\end{cases}
$$
且每次索赔发生的概率均为$p$。索赔的数量为

\begin{equation}
N=\sum_{i=1}^{n}X_i.
\label{eq:ch15-number-of-claims}
\end{equation}

其期望值为

\begin{equation}
\mathbb{E}[N]=np,
\label{eq:ch15-expected-claims}
\end{equation}
此外，在独立性方面，

\begin{equation}
\operatorname{Var}(N)=np(1-p).
\label{eq:ch15-claims-variance}
\end{equation}

观察到的索赔频率定义为

\begin{equation}
\hat p=\frac{N}{n}
=\frac{1}{n}\sum_{i=1}^{n}X_i.
\end{equation}

在独立分布下，其方差为

\begin{equation}
\operatorname{Var}(\hat p)
=\frac{1}{n^2}\operatorname{Var}(N)=\frac{p(1-p)}{n}.
\label{eq:ch15-frequency-variance}
\end{equation}

随着风险池扩大，索赔频率的方差按投保人数量的倒数下降。总索赔数量仍不确定，但平均每张保单的赔付经验变得更可预测。这构成风险汇集的统计基础，使保险公司能够更可靠地估计组合损失。

可预测性的提高取决于各种损失之间的关联性。当损失之间相关性较弱时，集中处理会更加有效；而当出现共同冲击时，可能会同时引发许多索赔，从而削弱分散风险的效果。

\begin{example}{汇集独立索赔}{pooling-independent-claims}

假设有10,000份类似的保单，每份保单的保额均为20,000人民币，且每年的索赔概率为1\%。

索赔的预期数量是

$$
10{,}000\times 0.01=100.
$$

索赔的标准差为

$$
\sqrt{10{,}000\times 0.01\times 0.99}
\approx 9.95.
$$

预期的索赔频率为1\%，而其标准差为

$$
\frac{9.95}{10{,}000}
\approx 0.001,
$$

或者大约0.10个百分点。

任何一位投保人的损失仍然不确定，而整个群体的索赔频率则变得更加稳定。

\end{example}


\subsection{可保性与风险汇集的限度}

更大保险池的多样化效益取决于各次损失之间的依赖性。为了说明这一点，设$X_i$表示第$i$份保单的损失。假设各次损失的方差均为$\sigma^2$，且每对损失都具有共同的相关性$\rho$。那么平均损失的方差就是

\begin{equation}
\operatorname{Var}(\bar{X})=
\frac{\sigma^2}{n}
+
\frac{n-1}{n}\rho\sigma^2.
\label{eq:ch15-correlated-pooling}
\end{equation}

当~$\rho=0$~时，随着风险池扩大，平均损失方差趋于零。当~$\rho>0$~时，第一部分随风险池规模下降，但相关性成分保留。因此，风险汇集可以减少个体特有波动，却不能消除共同冲击的风险暴露。

这种区分有助于解释为什么有些风险比其他风险更容易被承保。孤立的家庭火灾和许多汽车事故在保单的具体规定上存在很大差异，通常可以将这些事故合并起来处理。而地震、地区性洪水、流行病以及系统性网络事件则可能会同时导致许多投保人遭受损失，这就使得总风险更加难以分散。

因此，可保性不只取决于预期损失。事件与承保损失须足够明确且可核实，损失频率和程度须能够较可靠地估计，潜在总赔付也须与保险费及可用风险承担能力相匹配。信息和激励问题还可能限制可持续条款下能够提供的保障。

\begin{tcolorbox}[
enhanced,
breakable,
colback=white,
colframe=FinanceBlue,
boxrule=0.7pt,
sharp corners,
borderline west={2.2pt}{0pt}{FinanceBlue},
title=\textbf{财务展望：当资源整合仍不足以解决问题时},
coltitle=white,
colbacktitle=FinanceBlue,
fonttitle=\bfseries,
before skip=1.2em,
after skip=1.2em
]

当大部分损失风险由共同因素引起时，保险公司必须决定保留多少风险，以及将多少风险转移出去。地理上的多样化可以降低对局部事件的暴露风险，而再保险和与保险相关的证券可以将部分风险转移到其他资产负债表上。资本可以为那些仍由保险公司承担的损失提供缓冲。

这些机制并不能消除潜在的风险。它们只是将风险的财务后果重新分配给了不同的承担者，并决定了保险公司能够承受多大的集中性损失。

\end{tcolorbox}


\section{保险契约与信息问题}

保险契约规定了哪些损失是涵盖在内的，哪些损失由被保险人承担，以及索赔是如何被核实和支付的。这些条款还有助于管理在保险开始之前和之后出现的各种信息和激励问题。


\subsection{保险保障、契约支付与索赔}

保险单规定保险事件、保障期间、给付或赔偿、除外责任、免赔额、限额及索赔条件。这些条款共同决定何时触发支付，以及保险公司支付多少。

对于涵盖实际损失$L$的赔偿保单，令$d$表示免赔额，$C$表示保险公司的最大支付金额。简化的支付函数如下：

\begin{equation}
I(L)=\min\left\{\max(L-d,0), C\right\}.
\label{eq:ch15-insurance-payoff}
\end{equation}

被保险人所保留的损失份额为

\begin{equation}
R(L)=L-I(L).
\label{eq:ch15-retained-loss}
\end{equation}

因此，契约重新分配财务损失，而非消除损失。免赔额使部分损失由被保险人保留，赔付限额则限制保险公司的最大风险暴露。

这种支付结构对于赔偿保险来说最为直接。人寿保险通常在特定事件发生时提供一定的利益，而不是对实际财产损失的赔偿。健康保险则可能结合赔偿、合同利益、医疗服务提供者安排以及费用分摊等方式。因此，支付的形态取决于被保险事件的性质，以及可以观察和验证的内容。

索赔管理负责执行合同中的支付条款。保险公司会确定某项索赔是否属于保险范围，核实相关的损失或事件，按照保单条款进行赔付，并作出相应的支付。这些核实和结算程序也会影响到与合同相关的信息和激励措施。

\begin{example}{免赔额与保障限额}{deductible-and-limit}

该保单的免赔额为人民币~20,000，保险公司的最高赔偿金额为人民币~300,000。

如果承保损失为人民币80,000元，保险公司将予以赔偿

$$
I(80{,}000)=80{,}000-20{,}000=60{,}000.
$$

如果承保损失达到500,000元人民币，保险公司将予以赔偿

$$
I(500{,}000)=300{,}000,
$$

导致被保险人损失人民币约20万。

免赔额决定了保险公司需要承担多少较小的损失，而赔偿限额则限制了保险公司对较大损失的承受能力。

\end{example}


\subsection{逆向选择、核保与风险分类}

潜在投保人掌握自身风险信息，而保险公司无法完全观察时，可能产生逆向选择。如果不同风险类型面对相同条款，高风险者可能更愿购买或维持保障，低风险者则可能减少参与，导致承保风险池比定价时预计的更高风险。

假设有两种风险类型。类型$L$的预期损失为$\ell_L$，类型$H$的预期损失为$\ell_H$，

\begin{equation}
\ell_H>\ell_L.
\label{eq}
\end{equation}

基于两类风险共同确定的混合保险费，相对于~$L$~类的预期损失通常偏高，相对于~$H$~类则偏低。如果参与结构向~$H$~类转移，承保风险池的平均预期损失就会上升，对保险费形成上行压力，或要求调整契约条款。

核保通过收集和评估信息来解决这一信息问题，从而决定是否以及以何种条件为某项风险进行保险。根据产品和法律环境的不同，保险公司可能会考虑损失历史、财产特性、健康信息、商业活动、地理位置等因素，以及其他被允许的指标。

风险分类利用这些信息将具有相似预期索赔特征的保险风险进行分组。更详细的分类可以提高保险费与合同条款与预期损失之间的匹配度，不过其有效性取决于基础信息的质量、相关性以及允许使用的程度。

核保与风险分类通过改善承保风险评估及定价所依赖的信息，帮助保险公司应对逆向选择。相关评估主要关注损失的预期发生频率与程度，必要时还包括不同风险暴露之间的依赖关系。

\subsection{道德风险、风险分担与激励设计}

保障开始后，保险改变行为时可能产生道德风险。保障减轻被保险人承担的财务后果，可能削弱其预防损失、维护投保资产、控制损失程度或控制承保支出的激励，尤其在这些行动有成本且无法完全观察时。

保险契约可以让保险公司与投保人分担损失，而非完全转移损失。免赔额为~$d$~时，被保险人承担承保损失的初始部分。在比例分担安排下，被保险人还承担超过免赔额部分的一定比例损失。若保险公司承担免赔后损失的比例为~$\alpha$，则简化支付为


\begin{equation}
I(L)
=\alpha\max(L-d,0),
\qquad 0\leq\alpha\leq 1,
\label{eq:ch15-coinsurance-payoff}
\end{equation}
遵守任何覆盖限制规定。保留LaTeX命令、公式、交叉引用和数字。

风险分担使得被保险人面临更大的损失风险，这反而能够增强预防和成本控制的动力。更全面的保险覆盖则能提供更强的抵御不利状况的保护。因此，合同的设计需要平衡风险分担与激励效应之间的关系。

监控和索赔核实可以补充这种风险分担的形式。保险公司可能会要求采取预防措施，检查被保财产，核实损失和文件资料，使用经过批准的供应商网络，或者以特定的行为作为保险覆盖的条件。这样的机制可以提高激励效果并降低索赔成本，但其有效性取决于监控的可见性和监控成本。

\begin{tcolorbox}[
enhanced,
breakable,
colback=white,
colframe=FinanceBlue,
boxrule=0.7pt,
sharp corners,
borderline west={2.2pt}{0pt}{FinanceBlue},
title=\textbf{财务前景：报道的广度可能会削弱激励措施的效果。},
coltitle=white,
colbacktitle=FinanceBlue,
fonttitle=\bfseries,
before skip=1.2em,
after skip=1.2em
]

完全转移风险能提供最强的损失保护，但也可能降低被保险人采取行动的财务后果。保留部分损失可以维持激励措施，但代价是让投保人面临更大的不良后果风险。

因此，风险分担的适当程度取决于风险的性质、行为的可观察性、监控的成本，以及投保人承担剩余损失的能力。

\end{tcolorbox}


\section{保险定价与负债管理}

保险定价决定了为应对不确定的未来索赔及相关成本所需的保险费。这些索赔会产生负债，这些负债可能会持续到收取保险费之后，因此保险公司需要随着时间来管理储备金、资产、流动性、资本以及风险转移问题。


\subsection{预期损失、保险费与风险附加}

保险定价始于未来索赔支付的预期价值。用$L$表示与一份保单相关的保险损失的实际现值。预期损失是

\begin{equation}
EL=\mathbb{E}[L].
\label{eq:ch15-expected-loss}
\end{equation}

简化的保险费可以表示为

\begin{equation}
P=EL+E+R,
\label{eq:ch15-premium-decomposition}
\end{equation}

其中，$E$表示预期的获取和运营费用，而$R$则代表了不确定性的风险负荷，以及为了应对未被预期索赔所涵盖的风险而需要的资本。

预期损失既反映了索赔的频率，也反映了索赔的严重程度。在一个简单的设定中，该时期内最多只发生一次索赔。如果以概率$p$发生索赔，那么在已知有索赔发生的条件下，预期损失为$\mathbb{E}[S\mid \text{claim}]$。

\begin{equation}
EL=
p,\mathbb{E}[S\mid \text{索赔}].
\label{eq:ch15-frequency-severity}
\end{equation}

风险附加反映承担预期损失周围不确定性的成本。索赔波动较大、相关性较高、难以分散，或存在显著模型不确定性时，其重要性通常更大。再保险、监管、竞争、投资条件和契约设计，也可能影响最终收取的保险费。

因此，定价是面向未来的。历史损失数据可以提供信息，但当索赔频率、严重性、投入成本、法律条件、灾难风险或保单条款发生变化时，保险费可能需要进行调整。

\begin{example}{从预期损失到保险费}{expected-loss-to-premium}

保险公司估计每年的索赔概率为4\%，并且在发生索赔的情况下，预期索赔损失为约50,000人民币。每保单的预期损失为

$$
EL=0.04(50{,}000)=2{,}000.
$$

假设预期的收购和运营费用为500元人民币，而保险公司适用的风险系数为300。那么所确定的保险费为

$$
P=2{,}000+500+300=2{,}800.
$$

预期损失与所标示的保险费之间的差额，反映了在这种简化定价方式下所需的费用和风险负荷。

\end{example}


\subsection{保险负债估值与准备金}

定价将未来索赔的估算转化为当前的保险费条款。一旦保险合约签订，那些未来的索赔就变成了负债，其金额和支付时间必须在合同的整个期限内进行衡量。

让 $C_t$ 表示未来日期 $t=1,\ldots,T$ 时不确定的索赔现金流。作为一种简化的估值基准，预期索赔支付的现值可以表示为

\begin{equation}
V_L=\sum_{t=1}^{T}
\frac{\mathbb{E}[C_t]}{(1+r_t)^t},
\label{eq:ch15-insurance-liability-pv}
\end{equation}

其中，$r_t$表示与到期日和估值基准相适应的贴现率。该表达式说明了预期现金流和贴现的作用；更全面的负债衡量标准还可能反映费用、合同特性、不确定性以及其他相关估值框架所需的调整因素。

根据适用的会计或监管框架，储备金被表示为当前的负债。根据合同的不同，这些储备金可能反映那些已经发生但尚未支付或未报告的索赔，以及与长期保险契约相关的未来收益。例如，人寿和年金负债可能会受到死亡率、寿命、失效行为、费用、担保以及贴现率的影响。

因此，储备金并非仅仅是为未来索赔而预留的现金。它是一项由保险公司资产和资本支持的负债。索赔情况、假设或贴现率的变化可能会改变所计算的负债以及保险公司的财务状况。

预期支付的时间也很重要。与短期负债相比，长期负债通常对贴现率的变化更为敏感，因此将保险负债的估值与更广泛的资金时间价值框架联系起来是必要的。



\subsection{资产负债管理与偿付能力}

估值和预留将未来的保险义务转化为具有特定时间、敏感性和不确定性的负债。然后，保险公司管理资产以抵消这些负债，并保持资本充足，以应对资产负债表两边的逆向情况。因此，资产与负债管理着重于资产与负债现金流、流动性、久期以及风险的匹配。

设~$A$~表示资产经济价值，$L$~表示保险及其他负债的经济价值。经济盈余的一种简化度量为

\begin{equation}
S=A-L.
\label{eq:ch15-economic-surplus}
\end{equation}

资产负债表任一侧变化，都可能影响经济盈余。利率变化是重要来源之一。利率下降通常提高长期负债价值，具有类似久期特征的资产则可抵消部分风险暴露。

简化的久期-缺口衡量方法为

\begin{equation}
DG=D_A-\frac{L}{A}D_L,
\label{eq:ch15-insurer-duration-gap}
\end{equation}

其中，$D_A$~和~$D_L$~分别表示资产和负债的有效久期。久期缺口为经济盈余对利率变化的敏感性提供一阶度量。当保险负债包含保证、退保选择权、失效行为或其他非线性特征时，实际风险暴露可能更复杂。

流动性为资产-负债管理增加了另一个维度。财产保险公司可能在灾难发生后面临集中性的索赔支付，而人寿保险公司则可能会遇到资金退保或福利支付的问题。即使保险公司拥有足够的长期资产价值，但如果现金流出发生在资产能够以合理的成本转化为资金之前，那么该公司仍可能面临短期压力。

偿付能力涉及保险公司吸收资产和负债不利变化的能力。资本为应对意外索赔、资产损失、模型误差以及其他不利情况提供损失吸收能力。不同司法辖区的监管资本要求各不相同，但其基本经济作用是在实际结果不如预期负债和资产价值所反映的情况时提供缓冲。


\begin{figure}[H]
\centering
\begin{tikzpicture}[
> =Latex,
box/.style={
 draw,
 rounded corners,
 align=center,
 text width=5.0cm,
 minimum height=0.95cm,
 inner sep=4pt
 },
 main/.style={
 draw=FinanceBlue,
 fill=FinanceBlue!8,
 rounded corners,
 align=center,
 text width=5.0cm,
 minimum height=1.0cm,
 inner sep=4pt
 },
 flow/.style={-Latex, line width=0.8pt, FinanceBlue},
 link/.style={<->, line width=0.8pt, FinanceBlue}
 ]
\node[box] (assets) at (-3.4,1.6)
{资产\\流动性、久期、信用质量};

\node[box] (liab) at (3.4,1.6)
{负债\\时间、金额、敏感性};

\node[main] (alm) at (0,-0.5)
{资产--负债 \\ 管理};

\node[box] (capital) at (0,-2.2)
{保留LaTeX命令、公式、交叉引用和数字。  
大写\\损失吸收能力};

\draw[link] (assets) -- (liab);
\draw[flow] (assets) -- (alm);
\draw[flow] (liab) -- (alm);
\draw[flow] (capital) -- (alm);

\end{tikzpicture}

\caption{资产负债管理与损失吸收能力}
\label{fig}
\end{figure}


\section{保险机构与风险转移}

保险公司通过保险契约来承担风险，但它们并不需要将所有风险都保留在自己的资产负债表上。根据风险的性质和集中度，保险公司可以选择保留部分风险，将其中的一部分转移给再保险公司，或者将某些风险转移给资本市场投资者。

\subsection{原保险公司与主要险种}

在保险体系中，主要保险公司直接从家庭、企业以及其他投保人那里承担风险，涵盖人寿、健康以及财产保险等主要领域。

人寿保险提供与死亡或生存相关的福利。某些合同覆盖较长的期限，可能包含储蓄成分、担保或投保人的选择，这使得负债容易受到死亡率、寿命长短、退保行为和利率的影响。

健康保险涵盖医疗费用或特定的健康相关福利。索赔可能频繁发生，而使用情况可能会因保险覆盖而有所变化，因此在许多健康保险安排中，信息管理和激励措施显得非常重要。

财产保险保单涵盖财产损失、负债、机动车保险、商业中断险以及其他非人寿风险。索赔结算的时间跨度可以从较短的时段到多年不等，而灾难性事件可能导致各保单之间产生巨大的相关损失。

这些差异导致不同的负债状况，从而影响保险费设定、储备期限、流动性需求、资产与负债管理以及再保险的需求。



\begin{table}[H]

\centering

\caption{主要险种与负债特征}

\label{tab}

\begin{tabularx}{0.92\textwidth}{@{}p{0.20\textwidth} X X@{}}

\toprule

\textbf{行号} & \textbf{典型的保险风险敞口} & \textbf{负债特征} \\

\midrule

人寿保险
& 死亡率、生存率和长期家庭保障
& 通常具有较长的久期；可能包含担保和投保人的选择选项 \\

健康保险
& 医疗费用及与健康相关的福利
& 索赔可能频繁发生；使用和信息方面的问题可能会很严重 \\

属性与伤亡情况&财产损失、负债、机动车损害以及业务中断&索赔久期因行而异；灾难性和负债风险可能导致大额或延迟的支付\\

\bottomrule

\end{tabularx}

\end{table}


\subsection{再保险与保险风险再分配}


上述负债特征差异，会产生不同的风险转移需求。再保险使原保险公司支付再保险费，将部分承保风险暴露转移给再保险公司，从而在机构间重新分配承保损失。


再保险可以限制对个别重大损失的承受能力，减少灾难事件带来的损失，稳定承保结果，并提高保险公司开展额外业务的能力。保护的形式取决于保险费与损失如何在主保险公司和再保险公司之间分配。

在比例再保险中，保险费与损失是根据约定的比例来分配的。如果再保险公司的份额为$\alpha$，那么分担的损失就是

\begin{equation}
C(L)=\alpha L.
\label{eq:ch15-proportional-reinsurance-payoff}
\end{equation}

在损失超赔再保险中，再保险公司负责支付超出保留金$K$的赔付金额，但最高支付额有限，为$M$：

\begin{equation}
C(L)=\min\left\{\max(L-K,0),,M\right\}.
\label{eq:ch15-excess-of-loss}
\end{equation}

因此，主要保险公司将承担高达$K$的损失，以及超过$K+M$的任何金额。合同改变了保险公司的损失承担方式，尤其是在不利的情况下。

\begin{example}{超额损失再保险}{excess-loss-reinsurance}

该保险公司购买了价值约1亿元人民币的灾难性保险再保险，并且提供的最大赔付金额为约2亿元人民币。

如果涉及到的灾难性损失达到8000万人民币，那么再保险商将不支付任何赔偿。

如果损失达到1.8亿元人民币，再保险商将进行赔付。

[
180-100=80
]

百万人民币。

如果损失达到4亿元人民币，再保险公司将支付最多2亿元人民币的赔偿。主要保险公司将获得最初的1亿元人民币赔偿，而超过再保险层部分的1亿元人民币则由再保险公司承担。

\end{example}

通过接受来自多个主要保险公司的再保险风险，再保险公司可以将不同公司、地区以及业务领域的风险进行整合。这种风险分散的效果取决于这些风险之间的相互依赖性，而高度相关的事件可能会导致大量风险集中在再保险层面。

再保险也形成对再保险公司本身的风险暴露。因此，保障价值不仅取决于契约条款，也取决于承保损失发生时，再保险公司履行义务的能力。


\subsection{保险连接证券与巨灾债券}

再保险是将保险风险转移给其他保险报表，而证券-linked证券则可以将这种转移扩展到资本市场投资者。灾难债券就是最典型的例子。

典型巨灾债券结构中，投资者本金存入抵押安排，投资者因承担指定事件风险获得补偿。风险期间内若未发生契约触发事件，本金到期返还；若触发条件满足，则依契约释放部分或全部本金，减轻发起人对承保事件的风险暴露。

设债券的本金为$F$。忽略票息支付，简化的终端本金支付为

\begin{equation}
X_T=
\begin{cases}
F, & \text{如果没有触发任何事件发生的话。},\\
(1-\theta)F, & \text{如果触发条件导致本金的减少 } \theta,
\end{cases}
\label{eq:ch15-cat-bond-payoff}
\end{equation}

其中，$0\leq\theta\leq1$ 的值由合约的触发条件和损失结构来决定。

触发因素可能基于赞助方的实际损失、行业损失、物理事件参数、模拟出的损失，或者这些指标的组合。触发因素的选择会影响基础风险的分配，以及结算的透明度和速度。

灾难债券拓宽了用于吸收集中保险损失的资本来源。当灾难风险较大或需要额外的承担风险能力时，灾难债券可以补充再保险的作用。投资者因此能够因承担具有不同驱动因素的风险而获得补偿，这种风险与传统的金融市场风险有所不同。

这种转移并非无成本之举。结构设计费用、不确定性建模、基础风险、法律安排以及所需的投资者回报都会影响该证券的经济价值。因此，保险-linked证券应该是对传统再保险的补充，而不是简单地替代它。

\begin{tcolorbox}[
enhanced,
breakable,
colback=white,
colframe=FinanceBlue,
boxrule=0.7pt,
sharp corners,
borderline west={2.2pt}{0pt}{FinanceBlue},
title=\textbf{财务展望：将风险转移扩展到保险资产负债表之外},
coltitle=white,
colbacktitle=FinanceBlue,
fonttitle=\bfseries,
before skip=1.2em,
after skip=1.2em
]

保险风险可以通过一系列资产负债表来传递：

$$
\text{投保人}
\rightarrow
\text{主要保险公司}
\rightarrow
\text{再保险人}
\rightarrow
\text{资本市场投资者}.
$$

每一步都会重新分配因特定损失而产生的财务后果，但同时并不会消除潜在的不确定性。

\end{tcolorbox}


%========================================================
% Summary
%========================================================

\TailModule{summary}{本章小结}

保险公司通过状态依存合同来转移不确定的损失。投保人支付一定的保险费以获得保护，以应对特定的损失，而保险公司则汇集这些风险，并保持足够的财务能力来处理索赔。

当个体损失之间的相关性足够弱时，风险汇集可以使整体损失变得更加可预测。然而，相关性和共同的冲击会限制这种好处，从而增加对资本和风险转移的需求。

保险契约在风险分担与激励措施之间实现平衡。核保和风险分类可以应对逆向选择问题，而免赔额、比例分担制、监控以及索赔核实则可以缓解道德风险。

定价需要考虑预期损失、费用以及承担不确定性的成本。储备措施用于衡量未来的义务，资产与负债管理将资产与负债进行协调，而资本则提供吸收损失的能力。

风险不必留在主要保险公司手中。再保险可以将风险转移给不同的保险公司，而与保险相关的证券则可以将某些风险转移给资本市场投资者。


%========================================================
% Key Terms
%========================================================

\TailModule{keyterms}{关键术语}

保险（Insurance）；风险汇集（risk pooling）；大数定律（law of large numbers）；可保性（insurability）；相关损失（correlated losses）；逆向选择（adverse selection）；核保（underwriting）；道德风险（moral hazard）；免赔额与比例分担（deductibles and coinsurance）；预期损失（expected loss）；风险附加（risk loading）；保险负债（insurance liability）；准备金（reserves）；资产负债管理（asset--liability management）；偿付能力（solvency）；再保险（reinsurance）；巨灾债券（catastrophe bonds）；保险连接证券（insurance-linked securities）。

%========================================================
% Financial AI Lab
%========================================================

\TailModule{aiprompt}{金融~AI~实验室}

\textbf{实验室：对保险定价及风险转移分析的审计分析}

请一个AI系统来分析一家简化的财产保险公司。提供保单数量、预期索赔频率、预期索赔严重程度、运营费用、拟定的保险费或风险负荷、免赔额，以及简化的总损失再保险层。

将分析过程分为四个阶段进行实施。

\begin{enumerate}

\item \textbf{请确认预期损失和保险费。}

确保索赔频率和条件性严重性能够一致地结合使用。区分每份保单的预期损失与总体预期损失，并检查各项费用以及任何风险负荷是如何影响所确定的保险费的。

\item \textbf{检查合同设计。}

研究免赔额如何改变保险公司的赔付金额以及投保人的自留损失。确定所提供的信息是否足以计算其对索赔预期的影响；仅使用平均损失可能无法准确反映扣除免赔额后的预期赔付金额。将任何关于道德风险的索赔与实际提供的假设区分开来。

\item \textbf{追踪再保险层。}

在附件点下方、覆盖层内以及其极限之上，应用几种 aggregate-loss 结果。重新计算每笔情况下的再保险公司支付金额和主要保险公司的保留损失。

\item \textbf{测试相关损耗场景。}

将独立索赔替换为一种灾难性状况，这种状况会同时导致许多保单的损失。解释为什么更大的风险池可以减少特异性差异，而不会消除对共同冲击的暴露风险。

\end{enumerate}

最终报告应区分算术结果以及关于损失分布、投保人行为、风险之间的依赖性以及合同解释方面的假设。


%========================================================
% Discussion Questions
%========================================================

\TailModule{questions}{讨论题}

\begin{enumerate}

\item 为什么一个风险厌恶的家庭愿意支付高于预期损失的保险费呢？

\item 大数法则如何支持风险汇集，为什么正相关会限制这种效应？

\item 免赔额和共同承担费用如何影响风险分担以及投保人的激励措施？

\item 为什么当投保人有关于自身风险的私人信息时，汇集的保险费会产生逆向选择压力呢？

\item 保险公司储备与保险公司资本之间的经济区别是什么？

\item 再保险公司和灾难债券是如何将保险风险从主要保险公司转移出去的？

\end{enumerate}


%========================================================
% Applied Problems
%========================================================

\TailModule{appprobs}{应用题}


\begin{enumerate}

\item \textbf{预期损失与保险需求。}

一个家庭面临每年有3\%的概率会损失约20万人民币。

\begin{enumerate}

\item 计算预期的年度损失。

\item 将预期损失与每年7,000人民币的保险费进行比较。

\item 解释为什么风险厌恶型家庭仍然可能倾向于选择这种保险契约。

\item 请指出除了风险厌恶之外，还有哪些因素可能会影响家庭的决策。

\end{enumerate}

\item \textbf{合并独立的索赔。}

一家保险公司负责处理25,000种类似的风险。每种风险每年的索赔概率均为2\%，且各种索赔的结果是相互独立的。

\begin{enumerate}

\item 计算索赔的预期数量。

\item 计算索赔数量的标准差。

\item 计算索赔频率的标准差。

\item 解释索赔结果之间的正相关性如何会影响汇集资金所带来的分散化效益。

\end{enumerate}

\item \textbf{免赔额和支付限额。}

该保单的免赔额为人民币~10,000元，保险公司的最高赔偿金额为人民币~150,000元。

\begin{enumerate}

\item 计算保险公司对8,000元人民币、60,000元人民币和250,000元人民币损失所承担的赔偿金额。

\item 计算每个案例中被保险人所承担的损失金额。

\item 解释免赔额如何影响投保人在首次损失中仅获得部分赔偿的激励措施。

\end{enumerate}

\item \textbf{保险费分解。}

保险公司估计每保单的预期损失为人民币约1800元，预期费用约为人民币450元，风险负荷为人民币约250元。

\begin{enumerate}

\item 根据方程~\ref{eq:ch15-premium-decomposition}计算所要求的保险费。

\item 计算所示保险费与预期损失之间的差额。

\item 将这笔保险费表示为所指示保险费的百分比。

\item 解释为什么保险费与预期损失之间的区别不应被解释为会计利润。

\end{enumerate}

\item \textbf{负债评估与利率。}

保险公司预计在三年后支付约2000万元人民币。采用4\%的恒定年贴现率进行计算。

\begin{enumerate}

\item 计算预期付款的现值。

\item 使用3\%的贴现率重新计算现值。

\item 请说明这一变革的方向。

\item 将这一结果与长期保险公司资产-负债管理问题联系起来。

\end{enumerate}

\item \textbf{超额损失再保险。}

主要保险公司承担总灾难损失中的首笔人民币约5000万元，并购买再保险，以获得额外高达人民币约1.5亿元的保障。

\begin{enumerate}

\item 计算再保险商对总损失的责任赔偿金额，具体为：3000万人民币、1亿人民币、1.8亿人民币，以及2.6亿人民币。

\item 计算每起案件中的主要保险公司所承担的损失。

\item 确定再保险层的连接点和耗尽点。

\item 解释合同是如何改变主要保险公司所保留的损失分配的上限的。

\end{enumerate}

\end{enumerate}


%========================================================
% Case Analysis
%========================================================

\TailModule{caseanalysis}{案例分析}

\textbf{案例：2021年河南洪水与保险风险分担}

2021年7月，河南省遭受了严重的洪水袭击，短时间内产生了大量的索赔申请。截至7月28日，该省的保险业已经收到了约412,300份索赔报告，初步估计的保险损失约为人民币98.04亿元。当时，已有约43,400份索赔得到了解决，赔偿金额总计人民币6.09亿元。

索赔涉及多条保险条款。其中，机动车保险相关的索赔约为226,400件，造成的损失估计为人民币约6.412亿元；而家庭财产和农业保险方面的索赔也相当可观。在同一地区和同一事件下出现的大量损失，为研究当多个风险同时受到影响时保险是如何运作的提供了有用的分析框架。

从四个角度来分析这一案例。首先，解释地理相关性如何影响风险汇集的多样化效益。其次，探讨保单限额、免责条款、免赔额以及保险价值如何导致实际损失与索赔损失之间的差异。第三，解释为什么在索赔的金额和时机仍然不确定时，保险公司可能需要流动性。第四，讨论再保险如何重新分配由主要保险公司保留的部分损失。

上述数据仅为初步结果，不应被视为最终的行业损失或索赔金额。

\emph{原始资料来源：2021年7月29日河南省防汛新闻发布会；2021年7月30日发布的官方信息，援引了河南省银行和保险监管局的声明。}


%========================================================
% Further Reading
%========================================================

\TailModule{furtherread}{延伸阅读}

\begin{itemize}

\item Arrow, Kenneth J. (1963),
``Uncertainty and the Welfare Economics of Medical Care（不确定性与医疗保健的福利经济学）,''
\emph{American Economic Review}, 53(5), 941--973.

\item Pauly, Mark V. (1968),
``The Economics of Moral Hazard: Comment（道德风险经济学：评论）,''
\emph{American Economic Review}, 58(3), 531--537.

\item Rothschild, Michael, and Joseph E. Stiglitz (1976),
``Equilibrium in Competitive Insurance Markets: An Essay on the Economics of Imperfect Information（竞争性保险市场中的均衡：不完全信息经济学研究）,''
\emph{The Quarterly Journal of Economics}, 90(4), 629--649.

\item Froot, Kenneth A. (2001),
``The Market for Catastrophe Risk: A Clinical Examination（巨灾风险市场：案例考察）,''
\emph{Journal of Financial Economics}, 60(2--3), 529--571.

\end{itemize}


\ChapterTail{summary,keyterms,aiprompt,questions,appprobs,caseanalysis,furtherread}
